\section{Cardinals and the Axiom of Choice}
\label{sec:cumulative-hierarchy-04}

\begin{definition}[Equipotence]
Sets \(A\) and \(B\) are \emph{equipotent}, written \(A\approx B\), iff
a bijection \(A\to B\) exists.  Thus a cardinal, in the sense recalled
from \cref{def:ordinal-cardinal}, is an ordinal not equipotent
with any smaller ordinal.
\end{definition}

In \(\ZF\), a set \(A\) has an ordinal cardinal \(\abs A\) iff \(A\)
is well-orderable.  Do not write \(\abs A\) for an arbitrary set before
well-orderability has been established.  Under the Axiom of Choice, every set is
well-orderable and therefore has a unique cardinal.

\begin{theorem}[Hartogs]
For every set \(A\), there is a least ordinal \(h(A)\) for which no injection
\(h(A)\to A\) exists.  This theorem is provable in \(\ZF\), without the Axiom
of Choice.
\end{theorem}

\paragraph{Proof sketch.}
The well-orders carried by subsets of \(A\) form a set.  Collapse each such
well-order to its unique ordinal and use replacement to collect the resulting
ordinals.  Their supremum plus one cannot inject into \(A\).  Leastness
follows from ordinal well-ordering.  The set-coding of well-orders by
subsets of \(A\) and the replacement instance are supplied in
\cite{Jech2003}.

\begin{definition}[Order vocabulary]
A \emph{partial order} on a set \(A\) is a binary relation on \(A\) that
is reflexive, transitive, and \emph{antisymmetric}: whenever it relates
\(x\) to \(y\) and \(y\) to \(x\), then \(x=y\).  A \emph{partially
ordered set} is a set with a fixed partial order.  A \emph{chain} is a
subset on which the order is total.  An \emph{upper bound} of a subset is
an element that is greater than or equal to every member of that subset.  A
\emph{maximal element} is an element strictly below no element; a
\emph{greatest member} is an element greater than or equal to every
element.  One well-order \emph{end-extends} another
iff it extends it and every added element is greater than every original
element.
\end{definition}

\begin{theorem}[Equivalent forms of the Axiom of Choice]
Over \(\ZF\), each of the following implies the others:
\begin{enumerate}
  \item the Axiom of Choice;
  \item every set can be well-ordered;
  \item every partially ordered set in which every chain has an upper bound
        has a maximal element; and
  \item every product of a set-indexed family of nonempty sets is nonempty.
\end{enumerate}
\end{theorem}

\paragraph{Proof sketch.}
The implication from the Axiom of Choice to the maximal-element principle is proved by
choosing, whenever possible, a strict extension and applying transfinite
recursion; Hartogs' theorem prevents the construction from continuing through
all ordinals.  The maximal-element principle applied to partial well-orders
ordered by end-extension yields a well-order of an arbitrary set.  A well-order
of the union of a set-indexed family of nonempty sets selects the least member
of each factor, giving a member of the product.  Finally, a member of every
such product is a choice function.  This records the implication cycle and its
use of replacement; the set-coding details are imported from
\cite[Chapter 5]{Jech2003}.

The four clauses of the preceding equivalence theorem are equivalent over
the stated base theory; they are not definitions of one another.  Any use
of one clause in a proof must either cite that theorem or assume the clause
directly.

\begin{definition}[Aleph hierarchy, cardinal exponentiation, CH, and GCH]
By transfinite recursion, \(\aleph_0\defeq\omega\);
\(\aleph_{\alpha+1}\) is the least cardinal greater than
\(\aleph_\alpha\), which exists by Hartogs' theorem; and
\(\aleph_\lambda\defeq\sup\set{\aleph_\beta:\beta<\lambda}\) at a
limit ordinal \(\lambda\).  Under the Axiom of Choice, define
\(2^{\kappa}\defeq\abs{\Pow(\kappa)}\) for every cardinal \(\kappa\).
The \emph{continuum hypothesis} is the sentence
\[
  \mathsf{CH}\defeq(2^{\aleph_0}=\aleph_1),
\]
and the \emph{generalized continuum hypothesis} is
\[
  \mathsf{GCH}\defeq\forall\alpha\,(2^{\aleph_\alpha}=\aleph_{\alpha+1}),
\]
each read as the evident \(\Lang_{\in}\)-formalization of the displayed
cardinal equality.
\end{definition}

\begin{warning}[The Axiom of Choice hidden in cardinal notation]
The assertion that any two cardinalities are comparable, the assertion that
every set has an ordinal cardinal of the same size, and the unrestricted use of
cardinal arithmetic all use consequences of the Axiom of Choice.  A proof in \(\ZF\) must not import
those principles through notation.
\end{warning}

\begin{exercise}
\label{ex:cumulative-hierarchy-01}
Prove that \(\alpha\subseteq\beta\) whenever \(\alpha\leq\beta\) are
ordinals.  Identify where transitivity is used, and explain why no
transfinite induction is required.
\end{exercise}

\begin{exercise}
\label{ex:cumulative-hierarchy-02}
Compute \(\rank(\set{\varnothing,\set{\varnothing}})\) directly from the
recursive rank equation and verify the result from the hierarchy definition.
\end{exercise}

\begin{exercise}
\label{ex:cumulative-hierarchy-03}
Explain why the Mostowski collapse can fail to be injective when extensionality
of \(R\) is removed.  Supply a finite counterexample.
\end{exercise}

